\section{Dataset Distillation} \label{sec:26-background-datasetdistillation}

\subsection{Language Models and \ac{NLP} Foundations}
\label{sec:26-background-datasetdistillation-nlpfoundations}

\begin{foldable}
  Track~B of this thesis operates in the natural language processing domain, where the scale problem introduced in \S\ref{sec:21-background-deeplearning-lm} takes a specific form: the datasets required to fine-tune or continually train large language models are themselves prohibitively large, and the cost of processing them grows proportionally with their size.
  This subsection provides the language modelling and task foundations needed to understand both the problem and the solution proposed in Chapter~\ref{ch:50-research03}.
\end{foldable}

\begin{foldable}
  \textbf{Tokenization and word representations.}
  Text cannot be fed directly to a neural network as raw characters; it must first be converted to a numerical representation.
  Tokenization splits a string into a vocabulary of sub-word units using algorithms such as
  Byte-Pair Encoding (BPE) or WordPiece, balancing vocabulary size against the frequency of unknown tokens.
  Each token is then mapped to a dense vector (word embedding) in $\mathbb{R}^d$ via a learnable look-up table.
  Pre-trained static embeddings such as GloVe learn a fixed representation per token from co-occurrence statistics, while contextual embeddings --- the standard in modern \ac{NLP} --- vary the token's representation based on its surrounding context.
\end{foldable}

\begin{foldable}
  \textbf{Tokenization and dataset distillation.}
  The choice of tokenization algorithm has direct consequences for any method that assigns per-example importance scores, as those used in importance-based dataset distillation.
  Because \ac{BPE} and WordPiece produce variable-length token sequences for the same surface text, the model's loss is computed over sequences of different lengths across examples.
  Gradient norms --- a natural proxy for importance --- therefore reflect both the informational content of an example and the sheer number of tokens it contributes to the loss sum.
  Without normalization, longer examples accumulate larger gradient norms by arithmetic necessity rather than by informativeness, introducing a length bias into the importance ranking.
  A second artefact arises from vocabulary frequency: rare or domain-specific words are decomposed into multiple sub-word units, each of which has been seen less often during pre-training and therefore carries higher per-token loss.
  An example rich in rare terminology will exhibit elevated loss not because it is more informative about the task boundary, but because its tokens are lower-frequency in the pre-training distribution --- a distributional mismatch between pre-training and fine-tuning corpora that inflates apparent importance.
  Both artefacts --- length bias and vocabulary-frequency bias --- must be accounted for when computing importance scores over a heterogeneous corpus.
  Chapter~\ref{ch:50-research03} addresses this by normalising gradient norms by sequence length, partially decoupling the importance estimate from these tokenization-induced confounds and improving distributional fidelity of the resulting surrogate corpus.
\end{foldable}

\begin{foldable}
  \textbf{Transformer architecture.}
  The Transformer~\cite{vaswani2017attention} block processes an input sequence $(\mathbf{x}_1, \ldots, \mathbf{x}_n)$ through alternating multi-head self-attention and position-wise feed-forward sublayers, each wrapped in a residual connection and layer normalization.
  Multi-head self-attention allows each token to attend to all other tokens in the sequence simultaneously, computing a weighted combination of value projections where the weights are derived from query-key dot products:
  \begin{equation}
    \text{Attention}(Q, K, V) = \text{softmax}\!\left(\frac{QK^\top}{\sqrt{d_k}}\right)V.
    \label{eq:attention}
  \end{equation}
  Running $h$ attention heads in parallel and concatenating their outputs enables the model to capture multiple distinct relational patterns across the sequence.
  Positional encodings are added to the token embeddings to inject sequence order, which is otherwise absent from the permutation-invariant attention mechanism.
\end{foldable}

\begin{foldable}
  \textbf{The fine-tuning paradigm and the dataset distillation problem.}
  Pre-trained models such as BERT encode general linguistic knowledge acquired from large unlabelled corpora spanning hundreds of millions of words.
  Fine-tuning adapts this knowledge to a specific task by continuing gradient-based optimization on a labelled dataset: all model parameters are updated, and a lightweight task head is added on top of the encoder.
  The computational cost of fine-tuning is dominated by the number of gradient steps, each of which requires a full forward and backward pass through the model.
  This count scales linearly with both dataset size and the number of training epochs.
  For a corpus of 100k examples, three training epochs, and a batch size of 32, fine-tuning requires approximately 9,375 gradient steps; at 110M parameters per step for BERT-base, the total floating-point work is substantial even before accounting for validation and checkpointing.
  Critically, this cost is not incurred once: it compounds across hyperparameter searches (learning rate, warmup schedule, weight decay), cross-validation splits for rigorous evaluation, and continual learning scenarios where the model must be re-adapted as new data arrives.
  Each such repetition processes the same corpus in full, treating every example as equally necessary despite empirical evidence that model performance is driven by a small fraction of high-information examples.
  Dataset distillation targets this inefficiency directly.
  By replacing the full corpus with a compact surrogate $\tilde{\mathcal{D}}$ that preserves the distributional structure relevant to the task, the number of required gradient steps is reduced by an order of magnitude while final task accuracy is maintained.
  The surrogate functions as a distilled representation of the original corpus --- one that achieves distributional fidelity at a fraction of the computational cost.
\end{foldable}

\begin{foldable}
  \textbf{BERT and pre-training.}
  BERT~\cite{devlin2019bert} (Bidirectional Encoder Representations from Transformers) adapts the Transformer encoder for general-purpose language understanding by pre-training on two objectives over large unlabelled corpora: masked language modelling (predicting randomly masked tokens from their context) and next-sentence prediction (classifying whether two sentences are consecutive).
  The resulting representations encode rich syntactic and semantic information transferable to downstream tasks via fine-tuning.
  BERT-base (12 layers, 110M parameters) is fine-tuned on task-specific labelled data by adding a lightweight classification head and training end-to-end for a small number of epochs.
\end{foldable}

\begin{foldable}
  \textbf{\ac{NLI} task definitions.}
  The experiments in Chapter~\ref{ch:50-research03} span two families of language tasks.
  \emph{Text classification} assigns a single label to a document: topic classification, sentiment analysis (positive/negative).
  \ac{NLI} operates on pairs of sentences --- a premise and a hypothesis --- and predicts one of three relations: entailment (the hypothesis follows from the premise), contradiction (they are mutually inconsistent), or neutral (neither).
  \ac{NLI} is generally harder than classification because it requires modelling cross-sentence semantic relationships rather than the distributional features of a single document; the interaction space is quadratic in sentence length.
\end{foldable}

\begin{foldable}
  \textbf{\ac{NLI} versus text classification for dataset distillation.}
  The structural difference between text classification and \ac{NLI} has direct implications for how difficult dataset distillation is on each task family.
  In text classification, the label is determined by the distributional properties of a single document: topic classifiers respond to keyword frequencies and thematic vocabulary, while sentiment classifiers respond to lexical polarity and intensity markers.
  A surrogate corpus must therefore capture the marginal distribution of documents across class-defining features --- a well-studied space that importance-based selection methods can approximate by targeting high-gradient examples near decision boundaries.
  \ac{NLI} imposes a strictly harder requirement.
  The label is a function of the \emph{relationship} between a premise sentence and a hypothesis sentence; neither sentence alone determines the label, and the same premise can participate in examples with all three relation types depending on the hypothesis paired with it.
  This means the surrogate must preserve the joint distribution over sentence-pair relationships rather than the marginal distribution over individual sentences.
  The effective feature space is higher-dimensional: it includes not only the content of each sentence separately but also the semantic conjunction of the two.
  Two further consequences follow for importance scoring.
  First, gradient norms computed from a single-document classification head --- a natural proxy for example importance in text classification --- do not straightforwardly reflect the informativeness of a sentence pair, because the gradient is sensitive to how both sentences interact through the cross-attention computation rather than to either sentence in isolation.
  Second, ensuring diversity in the surrogate requires covering the space of premise-hypothesis conjunctions, not merely the space of individual sentence topics or lengths.
  Chapter~\ref{ch:50-research03} evaluates the proposed method on both task families and reports results separately for classification and \ac{NLI} benchmarks, providing evidence on whether importance-based selection can handle this structural distinction.
\end{foldable}

\begin{foldable}
  \textbf{Why corpus scale motivates dataset distillation.}
  Fine-tuning BERT on a large task corpus (e.g., $>$100k examples) requires significant computation per training run, and the cost multiplies across hyperparameter searches, continual learning updates, and ablation studies.
  Dataset distillation addresses this by finding a compact surrogate corpus $\tilde{\mathcal{D}}$ that, when used for fine-tuning, matches the performance achieved on the full corpus.
  The gain compounds when the same corpus is used repeatedly --- a common pattern in NLP benchmarking.
\end{foldable}

\subsection{Problem Definition}
\label{sec:26-background-datasetdistillation-definition}

\begin{foldable}
  Dataset distillation~\cite{zhao2020dataset} (also called dataset condenzation) seeks a small surrogate dataset $\tilde{\mathcal{D}}$ with $|\tilde{\mathcal{D}}| = K \ll N = |\mathcal{D}|$ such that a model trained on $\tilde{\mathcal{D}}$ performs approximately as well as one trained on the full dataset $\mathcal{D}$:
  \begin{equation}
    \tilde{\mathcal{D}}^{*} = \arg\min_{|\tilde{\mathcal{D}}|=K}
    \mathbb{E}_{\theta \sim \mathcal{A}(\tilde{\mathcal{D}})}\!\left[\mathcal{L}(\theta; \mathcal{D})\right],
    \label{eq:dd-objective}
  \end{equation}
  where $\mathcal{A}(\tilde{\mathcal{D}})$ denotes the learning algorithm applied to $\tilde{\mathcal{D}}$.
  The compression ratio $N/K$ is typically large --- ranging from $10\times$ to $100\times$ in practice --- and the surrogate may consist of either selected real examples (coreset selection) or synthesised examples that do not appear in $\mathcal{D}$.
\end{foldable}

\begin{foldable}
  Dataset distillation and knowledge distillation are complementary but distinct. KD compresses a \emph{model}: given a trained teacher, it produces a smaller student with similar predictive behaviour.
  DD compresses a \emph{dataset}: given a full training corpus, it produces a smaller surrogate that induces similar model behaviour when used for training.
  The two problems share the distributional fidelity requirement --- both the student and the surrogate corpus must represent the teacher's or full dataset's distribution faithfully --- but they differ in their input (a trained model vs.\ a raw dataset) and output (a model vs.\ a corpus).
\end{foldable}

\begin{foldable}
  A complete DD method for the text setting should satisfy three desiderata.
  First, it should be \emph{knowledge-informed}: the selection or synthesis of $\tilde{\mathcal{D}}$ should reflect each example's contribution to model learning, not just its frequency or surface features.
  Second, it should ensure \emph{distributional coverage}: the surrogate should represent the full range of semantic variation in $\mathcal{D}$, not just its high-frequency modes.
  Third, for text, it should produce \emph{human-readable output}: unlike image distillation where synthetic pixels need not be interpretable, text examples must be fluent and coherent to be usable for fine-tuning and to support interpretability and privacy auditing.
\end{foldable}

\subsection{Image Dataset Distillation}
\label{sec:26-background-datasetdistillation-image}

\begin{foldable}
  Image dataset distillation methods synthesise a small set of images by optimising them to
  preserve some property of training on the full dataset.
  Three optimization targets have been explored.
\end{foldable}

\begin{foldable}
  \textbf{Gradient matching}~\cite{zhao2020dataset} (DC) synthesises $\tilde{\mathcal{D}}$ such that the model gradients computed on $\tilde{\mathcal{D}}$ match those computed on $\mathcal{D}$ at each training step:
  \begin{equation}
    \min_{\tilde{\mathcal{D}}} \sum_{t} \left\| \nabla_\theta \mathcal{L}(\theta_t;
    \tilde{\mathcal{D}}) - \nabla_\theta \mathcal{L}(\theta_t; \mathcal{D}) \right\|^2.
    \label{eq:dc-objective}
  \end{equation}
  Because the gradient matching loss is differentiable with respect to the pixel values of the synthesised images, the images can be optimised directly by backpropagation.
  The resulting synthetic images often appear unnatural --- they encode gradient information rather than visual realism --- but they reliably reproduce the training dynamics of the full dataset on simple benchmarks such as MNIST and CIFAR-10.
\end{foldable}

\begin{foldable}
  \textbf{Trajectory matching}~\cite{cazenavette2022dataset} (MTT) extends gradient matching to the full training trajectory: it optimises $\tilde{\mathcal{D}}$ such that the model trained on it for $T$ steps mimics the parameter trajectory of a model trained on $\mathcal{D}$ for a longer expert trajectory.
  Matching the entire trajectory rather than single-step gradients improves performance on larger and more complex datasets.
\end{foldable}

\begin{foldable}
  \textbf{Distribution matching}~\cite{zhao2023dataset} (DM) takes a softer approach: instead of matching gradients directly, it minimises the discrepancy between the feature-space distributions of real and synthetic data.
  Distribution matching is more computationally scalable than gradient or trajectory matching because it does not require unrolling the training computation graph.
\end{foldable}

\begin{foldable}
  All three methods share a fundamental incompatibility with the text domain.
  They synthesise examples by gradient-based optimization in the continuous pixel space of images.
  Text tokens are discrete: there is no well-defined gradient of a loss with respect to an integer token index.
  Straight-through estimators and continuous relaxations partially address this, but they introduce approximation errors that compound across the typically long token sequences of text examples, and the resulting ``soft'' token sequences are not human-readable.
\end{foldable}

\subsection{Text Dataset Distillation}
\label{sec:26-background-datasetdistillation-text}

\begin{foldable}
  The discrete and sequential nature of text requires a fundamentally different approach to dataset distillation.
  Prior methods for text DD fall into two families.
\end{foldable}

\begin{foldable}
  \textbf{Soft-embedding methods} apply continuous optimization in the embedding space rather than the discrete token space.
  SLDD~\cite{sucholutsky2021soft} optimises a set of soft-label, continuous embedding vectors that, when used as training inputs, preserve model performance.
  DDTC~\cite{li2021data} and DDAL~\cite{maekawa2023dataset} similarly learn embedding sequences optimised by gradient matching against a differentiable fine-tuning model.
  DiLM~\cite{maekawa2024dilm} combines gradient matching with $k$-centre subset selection, identifying a small coreset of real examples whose gradient signal is most representative of the full dataset.
  These methods achieve good compression ratios on simple classification tasks but share two limitations.
  The optimised embeddings are \emph{architecture-locked}: the soft embeddings are tailored to a specific model architecture and cannot be used to fine-tune a different model.
  More fundamentally, none of these methods produce natural language text --- the output is a set of continuous vectors that are uninterpretable and cannot be shared, audited, or used outside the training pipeline.
\end{foldable}

\begin{foldable}
  \textbf{Clustering-based methods} select real examples by clustering the full dataset in embedding space and taking cluster representatives.
  DaLLME~\cite{tao2024textual} embeds all training examples using a large language model and applies $k$-means clustering, selecting the example closest to each centroid.
  This approach is fast, architecture-agnostic, and produces human-readable output (real examples rather than soft embeddings).
  However, $k$-means assigns uniform weight to each cluster regardless of its contribution to model learning: rare but semantically important examples are not distinguished from frequent low-information ones.
  The selection does not formally minimise any measure of distributional divergence, and the per-sample importance of selected examples is not considered.
\end{foldable}

\begin{foldable}
  \textbf{LLM-based generation.}
  An alternative to selecting from $\mathcal{D}$ is to generate a synthetic candidate pool by fine-tuning a language model on $\mathcal{D}$ and sampling from it.
  This avoids the gradient-flow problem entirely: the generation process is autoregressive and produces fluent, human-readable text.
  Privacy is a natural byproduct: since the synthetic examples are generated rather than extracted from the original corpus, they do not directly expose training data.
  This generation paradigm forms the generation stage of the TAKE pipeline (\S\ref{sec:54-research03-methodology}).
\end{foldable}

\begin{foldable}
  However, generation alone does not solve the selection problem.
  A large synthetic candidate pool must still be distilled to a compact surrogate of size $K$, and the selection step must address two residual problems shared by all prior text DD methods.
\end{foldable}

\begin{foldable}
  \textbf{Residual problem 1 --- hard-sample bias.}
  Naive per-sample importance scores based on model gradients are dominated by hard or noisy examples at convergence.
  As the model fits the training distribution, its loss on clean, easy examples falls to near zero; gradient norms are therefore large only for samples the model still struggles with --- which are disproportionately mislabelled or atypical examples.
  Selecting high-scoring samples by single-checkpoint gradient norms therefore biases the surrogate toward noise.
\end{foldable}

\begin{foldable}
  \textbf{Residual problem 2 --- coverage gap.}
  Greedy or clustering-based selection does not formally minimise the distributional divergence between $\tilde{\mathcal{D}}$ and the full importance-reweighted corpus.
  Even with corrected per-sample scores, a greedy selector may concentrate on the highest-scoring examples and neglect lower-scoring but semantically distinct regions of the distribution.
\end{foldable}

\begin{foldable}
  \S\ref{sec:27-background-importanceselection} introduces the two technical tools --- trajectory-integrated importance scoring and optimal transport --- that Chapter~\ref{ch:50-research03} uses to address both problems.
\end{foldable}
